\section{Estimating $N_{\rm max}$}\label{app:nmax}

Section~\ref{sec:expected_stop_iteration} derives $N_{\rm goal}$, the expected minimum sample
size at which the precision condition $\omega_{\rm HDI} \le \omega_{\rm goal}$ is first met.
This gives a principled lower bound on the data required for any precision-based stopping
algorithm.
For DPitG, however, $N_{\rm goal}$ is not sufficient for budget planning: the algorithm may
continue beyond $N_{\rm goal}$ whenever the posterior has not yet settled on a decisive
conclusion.
Practitioners therefore need a principled upper budget $N_{\rm max}$ --- the sample size
beyond which data collection halts regardless of decisiveness.

This appendix derives analytical approximations for the percentiles of the DPitG stopping
distribution, providing a principled basis for selecting $N_{\rm max}$.
Under the same normal approximation used in Section~\ref{sec:expected_stop_iteration},
we show that the $p$-th stopping percentile $N_p$ can be expressed in closed form for the
Reject case and in near-closed form for the Accept case.
A key finding is that the worst-case $\theta_{\rm true}$ for budget planning is not
$\theta_{\rm true} = 0.5$ (as for $N_{\rm goal}$) but rather $\theta_{\rm true}$ near a ROPE
boundary, where the decision may remain inconclusive regardless of how many observations are
collected.

\subsection*{Two-stage structure of DPitG stopping}

DPitG requires both precision and decisiveness simultaneously.
Because the precision condition blocks any stop before approximately $N_{\rm goal}$,
the stopping time decomposes as
\begin{equation}\label{eq:dpitg_stopping_time}
    T_{\rm DPitG} \;=\; \min\bigl\{\,N \ge N_{\rm goal} : \text{Decision is conclusive}\,\bigr\}.
\end{equation}
For $N \ge N_{\rm goal}$ the HDI is sufficiently narrow, and the outstanding question is
whether it sits entirely inside or entirely outside the ROPE.
We approximate the cumulative distribution function (CDF) of $T_{\rm DPitG}$ by
\begin{equation}\label{eq:dpitg_cdf_approx}
    \Pr\bigl(T_{\rm DPitG} \le N\bigr)
    \;\approx\; \Pr\bigl(\text{Decision conclusive at }N\bigr),
    \qquad N \ge N_{\rm goal}.
\end{equation}
This \textit{sticky-decision} approximation rests on the observation that once the HDI lies
entirely inside or entirely outside the ROPE, it tends to remain so as the posterior
concentrates further.
The approximation is accurate when $\theta_{\rm true}$ is not close to a ROPE boundary;
near the boundary the HDI can oscillate between decisive and indecisive states, and the
approximation underestimates the tail.
We address this case separately in the final subsection.

\subsection*{Probability of a conclusive decision at sample size $N$}

Under the normal approximation, $\hat\theta_N \approx \mathcal{N}(\theta_{\rm true},\,V/N)$
after $N$ observations, where $V = \theta_{\rm true}(1-\theta_{\rm true})$ for the Bernoulli
case.
Setting $\sigma_N = \sqrt{V/N}$, the 95\% HDI is approximately
$\bigl[\hat\theta_N - z_*\sigma_N,\;\hat\theta_N + z_*\sigma_N\bigr]$.
Define the normalised distances from $\theta_{\rm true}$ to each ROPE boundary:
\begin{equation}\label{eq:normalised_distances}
    d_R(N) \;=\; \frac{\mathrm{ROPE}_{\rm max} - \theta_{\rm true}}{\sigma_N},
    \qquad
    d_L(N) \;=\; \frac{\theta_{\rm true} - \mathrm{ROPE}_{\rm min}}{\sigma_N}.
\end{equation}
Both quantities are positive when $\theta_{\rm true}$ lies strictly inside the ROPE;
$d_R < 0$ when $\theta_{\rm true}$ lies above the ROPE, and $d_L < 0$ when it lies below.
Note that both scale as $\sqrt{N}$, growing without bound as data accumulate when
$\theta_{\rm true}$ is in the interior of either region.

The probabilities of each conclusive outcome are:

\begin{enumerate}[label=(\alph*)]
    \item \textit{Accept null} (applicable when $d_R > 0$ and $d_L > 0$):
    \begin{equation}\label{eq:pr_accept}
        \Pr(\text{Accept at }N) \;=\; \Phi\!\bigl(d_R - z_*\bigr) - \Phi\!\bigl(-d_L + z_*\bigr).
    \end{equation}
    \item \textit{Reject null above ROPE} (applicable when $d_R < 0$, i.e.\
    $\theta_{\rm true} > \mathrm{ROPE}_{\rm max}$):
    \begin{equation}\label{eq:pr_reject_above}
        \Pr(\text{Reject at }N) \;=\; 1 - \Phi\!\bigl(d_R + z_*\bigr).
    \end{equation}
    \item \textit{Reject null below ROPE} (applicable when $d_L < 0$, i.e.\
    $\theta_{\rm true} < \mathrm{ROPE}_{\rm min}$):
    \begin{equation}\label{eq:pr_reject_below}
        \Pr(\text{Reject at }N) \;=\; \Phi\!\bigl(d_L + z_*\bigr).
    \end{equation}
\end{enumerate}

Here $\Phi$ denotes the standard normal CDF and $z_* \approx 1.96$ for 95\% HDI.
Each probability increases monotonically with $N$ when $\theta_{\rm true}$ is strictly inside
or outside the ROPE, and approaches 1 as $N \to \infty$ --- consistent with the posterior
eventually producing a decisive verdict for any $\theta_{\rm true}$ not on the boundary.

\subsection*{Percentile formulas}

The $p$-th stopping percentile of DPitG is approximately
\begin{equation}\label{eq:np_dpitg}
    N_p \;\approx\; \max\!\bigl(N_{\rm goal},\;N_p^{\rm decision}\bigr),
\end{equation}
where $N_p^{\rm decision}$ is the sample size at which the conclusive-decision probability
first reaches $p$, obtained by inverting the relevant equation above.
When $N_p^{\rm decision} \le N_{\rm goal}$, the precision condition is the binding constraint
and DPitG stops near $N_{\rm goal}$; when $N_p^{\rm decision} > N_{\rm goal}$, the decision
condition dominates and more data are required.

\paragraph{Reject case.}
For $\theta_{\rm true} > \mathrm{ROPE}_{\rm max}$ (by symmetry, an analogous result holds
for $\theta_{\rm true} < \mathrm{ROPE}_{\rm min}$), inverting Equation~\ref{eq:pr_reject_above}
and solving for $N$ yields the closed form
\begin{equation}\label{eq:np_reject}
    N_p^{\rm reject}
    \;=\;
    \frac{\bigl(z_* + z_p\bigr)^2\,V}{\bigl(\theta_{\rm true} - \mathrm{ROPE}_{\rm max}\bigr)^2},
\end{equation}
where $z_p = \Phi^{-1}(p)$ is the $p$-th standard normal quantile
(e.g., $z_{0.75} \approx 0.674$, $z_{0.95} \approx 1.645$).
The denominator $(\theta_{\rm true} - \mathrm{ROPE}_{\rm max})^2$ encodes the signal-to-noise
structure: the farther $\theta_{\rm true}$ lies outside the ROPE, the fewer observations are
needed to render the decision conclusive, and $N_p^{\rm reject}$ falls well below
$N_{\rm goal}$ for large effect sizes.

\paragraph{Accept case --- symmetric scenario.}
When $\theta_{\rm true} = \theta_{\rm null}$ is centred in the ROPE ($d_R = d_L$),
inverting Equation~\ref{eq:pr_accept} gives
\begin{equation}\label{eq:np_accept}
    N_p^{\rm accept}
    \;=\;
    \frac{4\bigl(z_* + z_p'\bigr)^2\,V}{\Delta_{\rm ROPE}^2},
\end{equation}
where $z_p' = \Phi^{-1}\!\bigl((1+p)/2\bigr)$ accounts for the two-sided nature of the
Accept condition (e.g., $z_{0.75}' \approx 0.674$, $z_{0.95}' \approx 1.96 = z_*$).
For general (asymmetric) $\theta_{\rm true}$ inside the ROPE, $N_p^{\rm accept}$ must be
obtained numerically from Equation~\ref{eq:pr_accept}.

\subsection*{Relationship between $N_p$ and $N_{\rm goal}$}

For the symmetric Accept scenario, dividing Equation~\ref{eq:np_accept} by
$N_{\rm goal} = 4z_*^2\,V\,/\,\omega_{\rm goal}^2$
(Equation~\ref{eq:expected_stop_iteration_clt}) gives
\begin{equation}\label{eq:np_ratio}
    \frac{N_p^{\rm accept}}{N_{\rm goal}}
    \;=\;
    \left(1 + \frac{z_p'}{z_*}\right)^{\!2}
    \!\left(\frac{\omega_{\rm goal}}{\Delta_{\rm ROPE}}\right)^{\!2}.
\end{equation}
Under the default parameters of Section~\ref{sec:experimental_design}
($\omega_{\rm goal} = 0.8\,\Delta_{\rm ROPE}$, $z_* \approx 1.96$):

\begin{enumerate}[label=(\alph*)]
    \item $N_{50\%} \approx 1.15\,N_{\rm goal}$: the median DPitG stop is only modestly beyond
    $N_{\rm goal}$, with the Accept condition met quickly once precision is achieved.
    \item $N_{95\%} \approx 2.56\,N_{\rm goal}$: for the symmetric Accept scenario,
    setting $N_{\rm max} \approx 3\,N_{\rm goal}$ provides a 95\% budget guarantee.
    For the default setup ($N_{\rm goal} \approx 600$ at $\theta_{\rm true} = 0.5$),
    this gives $N_{\rm max} \approx 1{,}536$, consistent with the value of $N_{\rm max} = 1{,}500$
    used in Section~\ref{sec:experimental_design}.
\end{enumerate}

More generally, from Equation~\ref{eq:np_ratio} with $p = 95\%$ (so $z_p' = z_*$):
\begin{equation}\label{eq:n95_general}
    N_{95\%}^{\rm accept} = \frac{16\,z_*^2\,V}{\Delta_{\rm ROPE}^2},
\end{equation}
independent of $\omega_{\rm goal}$.
The ratio $N_{95\%}^{\rm accept} / N_{\rm goal} = 4(\omega_{\rm goal}/\Delta_{\rm ROPE})^2$
decreases as $\omega_{\rm goal}$ becomes a smaller fraction of $\Delta_{\rm ROPE}$: a more
demanding precision goal makes it easier for the HDI to land inside the ROPE once precision
is achieved, so $N_{95\%}$ is closer to $N_{\rm goal}$.

\subsection*{The boundary case and implications for $N_{\rm max}$}

A fundamental asymmetry distinguishes $N_{\rm goal}$ and $N_p^{\rm DPitG}$:

\begin{enumerate}[label=(\alph*)]
    \item $N_{\rm goal}$ is maximised at $\theta_{\rm true} = 0.5$, the point of largest
    Bernoulli variance $V = \theta(1-\theta)$.
    \item $N_p^{\rm decision}$ diverges as $\theta_{\rm true}$ approaches a ROPE boundary,
    irrespective of $\omega_{\rm goal}$.
\end{enumerate}

To see why, set $\theta_{\rm true} = \mathrm{ROPE}_{\rm max}$ exactly, so $d_R = 0$.
From Equations~\ref{eq:pr_accept}--\ref{eq:pr_reject_above}, the total probability of a
conclusive decision satisfies
\begin{align}
    \Pr(\text{conclusive at }N)\big|_{d_R=0}
    &= \underbrace{\bigl[1 - \Phi(z_*)\bigr]}_{\text{Reject above}}
     + \underbrace{\bigl[\Phi(-z_*) - \Phi(-d_L + z_*)\bigr]}_{\text{Accept}} \notag \\
    &\;\to\; 2\bigl[1 - \Phi(z_*)\bigr] \;\approx\; 5\%
    \quad \text{as } N \to \infty,\label{eq:boundary_conclusive}
\end{align}
since $d_L \to \infty$ and $\Phi(-d_L + z_*) \to 0$.
The 5\% ceiling is not a finite-sample artefact: as $N$ grows, $\hat\theta_N$ concentrates
around $\mathrm{ROPE}_{\rm max}$, and the HDI shrinks to straddle the boundary with
probability $\approx 95\%$.
The conclusive fraction ($\approx 5\%$ split equally between Accept and Reject) arises from
the small probability that $\hat\theta_N$ wanders sufficiently far from $\theta_{\rm true}$
in a given experiment, a probability that does not vanish with $N$ because $\sigma_N$ shrinks
at the same rate as the excess required for decisiveness.\footnote{This 5\% figure coincides
with the HDI credibility level in the sense that it equals $1 - \mathrm{CI}$ fraction.
For a 94\% HDI the ceiling would be 6\%, and so on. Intuitively, a boundary $\theta_{\rm true}$
is precisely the scenario in which the HDI has equal probability of just clearing the boundary
on either side, and the HDI level sets how often this clearing occurs.}

A practitioner who sizes $N_{\rm max}$ using $N_{\rm goal}$ alone will therefore
underestimate the budget needed when $\theta_{\rm true}$ is near a ROPE boundary, because
$N_{\rm goal}$ depends on $V(\theta_{\rm true})$ and is maximised at $\theta = 0.5$, whereas
the boundary inconclusive problem is unrelated to variance and affects only values near
$\mathrm{ROPE}_{\rm min}$ or $\mathrm{ROPE}_{\rm max}$.

\subsection*{Practical guidance}

\begin{enumerate}[label=(\alph*)]
    \item \textit{When $\theta_{\rm true}$ is expected to lie clearly inside or outside the
    ROPE}, use Equations~\ref{eq:np_reject}--\ref{eq:np_accept} to compute $N_{95\%}$ and
    set $N_{\rm max}$ accordingly.
    For the symmetric Accept scenario with $\omega_{\rm goal} = 0.8\,\Delta_{\rm ROPE}$,
    the rule of thumb $N_{\rm max} \approx 3\,N_{\rm goal}$ provides a 95\% stopping
    guarantee.

    \item \textit{When $\theta_{\rm true}$ may lie near a ROPE boundary} --- signalling that
    the prior on the effect size is genuinely uncertain --- the appropriate question shifts
    from ``what $N_{\rm max}$ ensures 95\% stopping?'' to ``what inconclusive rate is
    acceptable at a given $N_{\rm max}$?''
    Equation~\ref{eq:boundary_conclusive} shows that near the boundary, the inconclusive rate
    is approximately $95\%$ regardless of $N_{\rm max}$.
    In such settings, $N_{\rm max}$ is a policy choice informed by the cost of
    an inconclusive outcome, not a statistical one.
    Reconsidering the ROPE width or the precision goal $\omega_{\rm goal}$ may be more
    productive than simply increasing the budget.
\end{enumerate}
